% Generality: how widely the same design choice is made, and what that does and does not show.

\section{How common is this design choice?}
\label{sec:generality}

One board is one board. The results above show that the verifier changes the ranking on
MATH-Hard, and a reader is entitled to ask whether that is a peculiarity of this benchmark or
of the frameworks that score it. The question splits into two, and only one of them is
answered here.

The answerable question is how often a scoring configuration coerces an unreadable response
into a wrong answer, which is the design choice that creates the problem in the first place.
We resolved the task configurations of four evaluation frameworks
\citep{biderman2024lessons, opencompass2026, fourrier2024lighteval, openai2024simpleevals} with each framework's own
loader at a pinned commit, and classified the scoring path each generative task actually
takes. In \texttt{lm-evaluation-harness}, \LmEvalCoerce{} of \LmEvalGenerative{} generative
tasks, or \LmEvalCoercePct\%, recover an answer by matching an assumed output format and
record a failure to match as an incorrect answer. In OpenCompass the figure is
\OpenCompassCoerce{} of \OpenCompassGenerative{}, in \texttt{lighteval}
\LightevalCoerce{} of \LightevalGenerative{} generative metrics, and in OpenAI's
\texttt{simple-evals} \SimpleEvalsCoerce{} of \SimpleEvalsStringMatch{} string-matching
modules.

Reading code is not the same as running it, so the classification was validated behaviourally.
For a random sample of tasks the declared filter pipeline was rebuilt from the framework's own
classes and run on a well-formed response and on an equivalent response carrying the same
answer in ordinary prose. Of \CensusValidTestable{} testable sampled tasks, all
\CensusValidConfirmed{} behaved as classified, and of \CensusControlN{} control tasks drawn
from the tasks classified as unaffected, all \CensusControlClean{} were confirmed unaffected.

The unanswerable question, with the evidence in this paper, is how often that design choice
changes a published ranking on some other benchmark. It depends on how often real responses
take a form the extractor misses, which varies by benchmark, by model, and by prompt, and we
measured it only here. The prevalence figures above bound how many tasks could be affected.
They do not establish that any particular one is, and we do not claim otherwise.

The strongest statement the evidence supports is narrower than the one a reader might expect,
and worth stating in full. On one board, with the leaderboard's own published generations, the
choice of verifier moves the ranking. The design choice that makes this possible is the
default across the major open evaluation frameworks. Whether it bites elsewhere is a
measurement someone has to make, and the census says where to look first.
